% Author-written (rendered from next.tex).
\section{Individual Fairness in Clustering}
Individual fairness in clustering aims to ensure that similar individuals are treated similarly. This concept was introduced by Jung, Kannan, and Lutz~\cite{JungKL20}, who formalized the \emph{individually fair $k$-clustering} problem by requiring each point to be assigned to a center within a radius equal to that of its $\lceil n/k \rceil$-th nearest neighbor. This fairness radius $\delta(x)$ ensures that no point is excessively far from a center compared to what is typical in a fair division of the dataset.

Their work provided a $2$-approximation for this fairness constraint but lacked control over the clustering cost. Follow-up work by Mahabadi and Vakilian~\cite{MahabadiV20} generalized the model to include $\ell_p$ clustering objectives, notably $k$-median and $k$-means, under individual fairness constraints. They designed a local-search algorithm with $(84, 7)$-bicriteria approximation, where the first parameter bounds the cost and the second bounds fairness violation.

Subsequent work improved both theoretical and practical aspects. Chakrabarty and Negahbani~\cite{ChakrabartyN21} used LP rounding to obtain a $(2p+2, 8)$-bicriteria approximation, giving improved guarantees especially for $p = 1$ (i.e., $k$-median). Vakilian and Yalçiner~\cite{VakilianY22} introduced a stronger LP formulation, achieving $(16p+\varepsilon, 3)$ approximation for general $p$, and $(7.08 + \varepsilon, 3)$ for $k$-median, with significant reductions in fairness violations.

Recently, Bateni et al.~\cite{bateni2024scalable} presented a scalable local-search algorithm for individually fair $k$-means clustering. Their method runs in $\widetilde{O}(nd + nk^2)$ time and achieves $(O(1), 6)$-approximation, showing superior scalability and empirical performance over previous methods. Notably, they introduced the concept of \emph{anchor zones} to ensure fairness constraints during local search swaps, enabling efficient enforcement of fairness while minimizing cost. Their algorithm was able to scale to large datasets with hundreds of thousands of points, outperforming prior LP- and SDP-based methods in both speed and empirical cost.

Ebbens et al.~\cite{ebbens2024subquadratic} further improved the theoretical landscape for the fair $k$-center variant by presenting a $(2,2)$-bicriteria approximation in $O(n^2 + kn\log n)$ time and a randomized $(10, 2 + \varepsilon)$ approximation in $O(nk \log(n/\delta) + k^2/\varepsilon)$ time, based on a novel fair radius estimation method. Their results match the best known approximability under standard complexity assumptions.

Overall, individually fair clustering has matured into a robust framework with both theoretical guarantees and scalable algorithms across different clustering objectives and norms, making it a practical and principled fairness constraint in unsupervised learning.
